\begin{center}\bfseries 9.\end{center}
\noindent\hspace*{\parindent}Nunmehr aber beginnt das üppig fruchtbare Neuland, um welches \textsc{Gauss}’ eigenste Forschung die Lehre von den quadratischen Formen bereichert hat. Zur Einteilung der Formen mit gegebener Determinante \(D\) in Klassen tritt
\newpage\SourceLeaf{29}\ArithmeticHead{22}{L}{BACHMANN, ÜBER GAUSS’ ZAHLENTHEORETISCHE ARBEITEN.}
\noindent hinzu die weitere Einteilung dieser Klassen in Ordnungen, je nach dem grössten gemeinsamen Teiler \(m\) der Koeffizienten \(a,\,2b,\,c\) ihrer Formen \((a,b,c)\), der für alle Formen einer Klasse derselbe ist. Hervorzuheben sind hier die eigentlich — und eventuell die uneigentlich — primitive Ordnung, für welche \(m=1\) resp. \(m=2\) ist. Die Klassen einer primitiven Ordnung zerfallen nun wieder in mehrere Geschlechter. Für alle zu \(2D\) teilerfremden Zahlen \(n\) nämlich, welche durch eine primitive Form darstellbar sind, sind ihre quadratischen Charaktere bezüglich der einzelnen Primfaktoren der Determinante und in besonderen Fällen auch bezüglich der Moduln 4 oder 8 die gleichen. Diese Einzelcharaktere, deren Anzahl \(\lambda\) heisse, bilden zusammen den Gesamtcharakter der Form und auch der Klasse, der sie angehört, und alle Klassen, deren Gesamtcharakter derselbe ist, bilden ein Geschlecht. Die Anzahl aller denkbaren Gesamtcharaktere ist \(\chi=2^\lambda\), und es entsteht die Frage, ob für jeden von ihnen ein entsprechendes Geschlecht wirklich vorhanden ist.

Zu ihrer Beantwortung dient eine eigentümliche Rechnung, deren Elemente nicht Zahlen, sondern Formenklassen sind. Sie gründet sich auf die sogenannte Zusammensetzung quadratischer Formen, für welche die oben (S. 18) eingeführte \textsc{Gauss}sche Grundformel das einfachste Beispiel ist: eine Form
\[
AX^2+2BXY+CY^2
\]
heisst aus den beiden andern:
\[
ax^2+2bxy+cy^2,\ a'x'^2+2b'x'y'+c'y'^2
\]
zusammengesetzt, wenn sie durch eine bilineare Substitution, deren Koeffizienten gewisse Bedingungen erfüllen, in ihr Produkt übergeht. Aus der von \textsc{Gauss} wieder im Gewande algebraischer Beziehungen gehaltenen ausführlichen allgemeinen Theorie solcher Zusammensetzung entnehmen wir hier nur die wesentlichsten Folgesätze. Ist \(F\) eine aus den Formen \(f,f'\) und \(F_1\) eine aus den Formen \(f_1,f'_1\) zusammengesetzte Form, so gehören \(F,\ F_1\) derselben Klasse an, wenn sowohl von \(f,f_1\) als auch von \(f',f'_1\) je dasselbe gilt. Somit darf die Klasse von \(F\) aus den Klassen von \(f,f'\) zusammengesetzt heissen.

\textsc{Gauss} fasst diese Zusammensetzung zweier Klassen \(C,\ C'\) als eine additive Verknüpfung \(C+C'\) derselben; zweckmässiger ist es, sie als eine Multipli-\newpage
\SourceLeaf{30}\ArithmeticHead{23}{R}{BINÄRE QUADRATISCHE FORMEN.}
\noindent kation zu fassen und die zusammengesetzte Klasse als Produkt der beiden andern durch \(C\mathbin{.}C'\) zu bezeichnen. Auf Grund dieser Definition liefert nun die Gesamtheit der primitiven Formenklassen — abgesehen etwa von den Potenzen einer primitiven Wurzel (\(\mathrm{mod.}\ p\)), auf deren Analogie \textsc{Gauss} selbst (art. 306, VI) hingewiesen hat — das erste Beispiel eines Begriffes, der bald nach \textsc{Gauss} in der gesamten Mathematik weitreichende Herrschaft eingenommen hat, des Begriffs der Gruppe, speziell der \textsc{Abel}schen Gruppe. Die Zusammensetzung der Formenklassen ist kommutativ und assoziativ; nennt man die Form \(x^2-Dy^2\) die Hauptform, die Klasse, der sie angehört, die Hauptklasse, das Geschlecht, welches die letztere enthält, das Hauptgeschlecht, so spielt die Hauptklasse bei der Zusammensetzung die Rolle der Einheit, indem jede Klasse bei der Zusammensetzung mit ihr ungeändert bleibt; zwei entgegengesetzte Klassen, nämlich solche, in denen entgegengesetzte Formen \((a,b,c),\ (a,-b,c)\) auftreten, sind einander reziprok d. h. sie setzen sich zur Hauptklasse zusammen.

Aus diesen von \textsc{Gauss} hergeleiteten Sätzen folgt nun nach einem allgemeinen für \textsc{Abel}sche Gruppen von \textsc{Kronecker}\footnote{Berliner Monatsberichte 1870, S. 881, Werke I, S. 271.}\ bewiesenen Satze, was speziell für die Gruppe der Formenklassen vorher schon \textsc{Schering}\begingroup\renewcommand{\thefootnote}{**)}\footnote{\textit{Die Fundamental-Classen der zusammensetzbaren arithmetischen Formen}, Göttinger Abhandlungen \textbf{14} (1869), Werke I, S. 135.}\endgroup\ gezeigt hat, dass alle Klassen eindeutig aus einer kleinsten Anzahl fundamentaler Klassen zusammensetzbar sind. Soweit ist aber \textsc{Gauss} noch nicht gelangt, wenigstens enthalten seine Aufzeichnungen nichts, was sich darauf bezieht; er hat nur noch an einer späteren Stelle der D. A. (in artt. 305 u. 306) sowie in einem nachgelassenen, aus dem Jahre 1801 (s. W. II, S. 268) stammenden Fragmente (W. II, S. 266) für diejenige Untergruppe, welche aus den Klassen des Hauptgeschlechtes besteht, solche Darstellung ihrer Klassen mittels fundamentaler begonnen. Lassen sie sich sämtlich durch den Zyklus \(C,\ C^2,\ C^3,\ .\,.\,.,\ C^k\) der verschiedenen Potenzen einer einzigen darstellen, so nennt er die Determinante \(D\) regulär, andernfalls irregulär, und bezeichnet, wenn dann jener Zyklus der grösste der vorhandenen ist, den Quotienten \(\frac{g}{k}\), in welchem \(g\) die Anzahl aller Klassen des Hauptgeschlechts bedeutet, als den Exponenten der Irregularität; der art. 306, VIII. (W. I, S. 374) enthält hierüber noch mancherlei Aussagen,
\newpage\SourceLeaf{31}\ArithmeticHead{24}{L}{BACHMANN, ÜBER GAUSS’ ZAHLENTHEORETISCHE ARBEITEN.}
\noindent u. a. dass die Determinante stets regulär ist, wenn \(g\) keinen quadratischen Teiler hat, dass in der wachsenden Reihe der negativen Determinanten die Menge der irregulären stets zuzunehmen scheint, usw.
